\chapter{Especificación de requisitos}\label{cap:especificacion}

Este capítulo presenta la especificación de requisitos de CloudPorter en dos niveles.
En primer lugar se describe la visión completa del sistema, el problema que CloudPorter
aspira a resolver en su totalidad, partiendo de una identificación clara de los
usuarios a los que va dirigida la herramienta. A continuación se delimita el alcance
del trabajo, distinguiendo los requisitos que se abordan en este TFG de aquellos
que quedan como líneas de evolución futura.

\section{Usuarios objetivo}

CloudPorter está dirigido a dos perfiles principales que comparten el problema del
cloud lock-in pero lo experimentan desde ángulos distintos.

El primero es el \textbf{desarrollador o estudiante} que se acerca a la nube
por primera vez o que trabaja en proyectos personales o de pequeña escala. Para
este perfil, la barrera principal es el conocimiento: aprender las tecnologías
específicas de un proveedor supone una inversión de tiempo significativa, y el
coste de comprometerse con uno sin poder comparar alternativas de forma sencilla
desincentiva la experimentación.

El segundo es el \textbf{equipo técnico o profesional} que afronta la elección de
proveedor al diseñar un nuevo despliegue, o que quiere conservar la libertad de
cambiar de proveedor en el futuro. Para este perfil, el problema es organizacional
y económico. La decisión implica estudios de costes comparativos, dependencias del
conocimiento del equipo y el riesgo de quedar atado a un proveedor si las condiciones
cambian.

\section{Visión del sistema}

CloudPorter se concibe, en su forma ideal, como una herramienta que abordaría el
cloud lock-in en todas sus dimensiones, desde la infraestructura hasta el código, y
desde la CLI hasta una plataforma de gestión completa. Esta sección describe esa
visión completa y aspiracional, sin restricciones de alcance; el subconjunto que
este trabajo aborda realmente se delimita en la sección siguiente.

\subsection{Núcleo de abstracción e infraestructura}

\begin{itemize}
    \item \textbf{Manifest cloud-agnóstico:} descripción de infraestructura en un
    formato YAML independiente del proveedor, que cubra cómputo, red, bases de
    datos, almacenamiento y otros recursos comunes.
    
    \item \textbf{Soporte multi-proveedor:} traducción del Manifest a multiples
    cloud provider: AWS, Azure, GCP, Oracle Cloud, Alibaba Cloud u otros.

    \item \textbf{Generación desde fuentes existentes:} capacidad de partir de
    un \texttt{docker-compose.yml} o una plantilla \textbf{Terraform} existente y
    traducirlos automáticamente a un CloudPorter Manifest listo para despliegue.

    \item \textbf{Cobertura de lock-in de código:} integración con herramientas
    de abstracción a nivel de SDK, como \textbf{MultiCloudJ}, para cubrir
    también servicios con lock-in de código, como almacenamiento de objetos
    o colas de mensajes.
\end{itemize}

\subsection{FinOps y costes}

\begin{itemize}
    \item \textbf{Estimación comparativa previa al despliegue:} cálculo del
    coste aproximado de una arquitectura en distintos proveedores antes de
    ejecutar el despliegue, con recomendaciones de optimización.

    \item \textbf{Monitorización del gasto:} visibilidad del gasto cloud en
    tiempo real, detección de recursos huérfanos y alertas de desviación
    presupuestaria.
\end{itemize}

\subsection{Seguridad}

\begin{itemize}
    \item \textbf{Análisis de la infraestructura generada:} integración con
    herramientas de análisis estático de IaC para detectar configuraciones
    inseguras —permisos excesivos, puertos expuestos— antes del despliegue.

    \item \textbf{Auditoría continua:} integración con herramientas de auditoría
    de seguridad cloud para la verificación continua del cumplimiento normativo.
\end{itemize}

\subsection{Interfaz y experiencia de usuario}

\begin{itemize}
    \item \textbf{CLI:} interfaz de línea de comandos como punto de entrada
    principal, accesible sin conocimiento previo de los proveedores.

    \item \textbf{Plataforma web:} interfaz gráfica para gestión visual de la
    infraestructura, con diagramas de arquitectura interactivos y paneles
    de monitorización, sistemas de métricas y alertas.

    \item \textbf{Integración con agentes externos:} exposición de las
    funcionalidades de CloudPorter como herramientas invocables por agentes
    de IA mediante el protocolo \hyperlink{mcp}{\textbf{MCP}} (\textit{Model Context Protocol})
    \cite{mcpDocs}, de forma que un modelo de lenguaje pueda operar la
    herramienta de forma autónoma.
\end{itemize}

\subsection{Componente agéntico}

\begin{itemize}
    \item \textbf{Asistente en lenguaje natural:} componente que permite
    interactuar con CloudPorter mediante lenguaje natural, asistiendo al
    usuario tanto en la generación del Manifest como en el uso general
    de la herramienta.

    \item \textbf{Base de conocimiento contextual:} sistema de recuperación
    aumentada de información (\textit{Retrieval-Augmented Generation}, \hyperlink{rag}{RAG}) que
    permite enriquecer al agente con documentación propia — diagramas, diseños
    técnicos, decisiones de arquitectura — para obtener sugerencias adaptadas
    al contexto real del proyecto.

    \item \textbf{Sugerencias de optimización:} recomendaciones de costes,
    arquitectura y seguridad generadas por el agente a partir del contexto
    del despliegue y la base de conocimiento.

    \item \textbf{Capa determinista final:} independientemente del componente
    de IA utilizado, toda infraestructura generada pasa por una capa de
    validación y traducción determinista antes del despliegue, garantizando
    resultados predecibles y auditables.

\end{itemize}

\section{Alcance de este trabajo}

El sistema descrito en la sección anterior representa la visión completa y ambiciosa de CloudPorter.
Dado el alcance de un Trabajo de Fin de Grado, se define un subconjunto acotado
de requisitos suficiente para demostrar la viabilidad del enfoque y aportar valor
real, sin comprometer la calidad de lo que se construye. Este subconjunto es el que
delimita la prueba de concepto que constituye el núcleo del trabajo: lo
mínimo necesario para evidenciar que el enfoque funciona de extremo a extremo.

\section{Requisitos funcionales}

Los siguientes requisitos funcionales definen el comportamiento del sistema en
su versión inicial. Se presentan agrupados por área funcional.

\subsection{Manifest y traducción}

\begin{itemize}
    \item \textbf{RF-01.} El sistema debe permitir definir infraestructura mediante
    un fichero YAML con un formato propio independiente del proveedor cloud.

    \item \textbf{RF-02.} El sistema debe validar el Manifest antes de procesarlo,
    informando de errores de sintaxis o valores no soportados.

    \item \textbf{RF-03.} El sistema debe traducir el Manifest a plantillas
    \textbf{OpenTofu} desplegables en \textbf{AWS} y \textbf{Microsoft Azure}.
    La elección de ambas plataformas y la preferencia por OpenTofu sobre
    Terraform se justifican en el capítulo~\ref{cap:planificacion}.

    \item \textbf{RF-04.} El sistema debe soportar los siguientes tipos de
    recursos en la traducción: instancias de cómputo, bases de datos relacionales
    y configuración de red —generada automáticamente, sin requerir declaración
    explícita del usuario.

    \item \textbf{RF-05.} El sistema debe resolver automáticamente las equivalencias
    entre recursos genéricos del Manifest y los servicios específicos de cada
    proveedor mediante un \hyperlink{mapper}{mapper} determinista.
\end{itemize}

\subsection{Despliegue}

\begin{itemize}
    \item \textbf{RF-06.} El sistema debe ejecutar el despliegue de la
    infraestructura traducida en el proveedor seleccionado por el usuario.

    \item \textbf{RF-07.} El sistema debe informar del estado del despliegue
    y de los recursos creados una vez completado.
\end{itemize}

\subsection{FinOps}

\begin{itemize}
    \item \textbf{RF-08.} El sistema debe estimar el coste mensual aproximado
    de desplegar la arquitectura definida en el Manifest para cada proveedor
    soportado.

    \item \textbf{RF-09.} El sistema debe presentar la comparativa de costes
    entre proveedores antes de ejecutar el despliegue, permitiendo al usuario
    tomar una decisión informada.
\end{itemize}

\subsection{Interfaz de línea de comandos}

\begin{itemize}
    \item \textbf{RF-10.} El sistema debe exponer su funcionalidad a través de
    una CLI que permita, como mínimo: validar un Manifest, estimar costes y
    ejecutar un despliegue.
\end{itemize}

\subsection{Seguridad}

\begin{itemize}
    \item \textbf{RF-11.} El sistema debe ofrecer, antes del despliegue, un análisis
    de seguridad de la infraestructura definida en el Manifest que detecte
    configuraciones o anti-patrones inseguros e informe al usuario.
\end{itemize}

\section{Requisitos no funcionales}

Más allá de las funcionalidades concretas, el sistema debe satisfacer una serie de propiedades
de calidad que condicionan su diseño y su evolución. Los siguientes requisitos no funcionales
las recogen:

\begin{itemize}
    \item \textbf{RNF-01. Determinismo:} dado un mismo Manifest y un mismo
    proveedor, el sistema debe generar siempre la misma plantilla de
    infraestructura, garantizando resultados reproducibles y auditables.

    \item \textbf{RNF-02. Extensibilidad:} la arquitectura del sistema debe
    permitir añadir soporte a nuevos proveedores cloud sin modificar el núcleo
    de la herramienta.

    \item \textbf{RNF-03. Usabilidad:} el sistema debe poder utilizarse sin
    conocimiento previo de los servicios específicos de AWS o Azure: el usuario
    opera en términos del Manifest, no del proveedor.

    \item \textbf{RNF-04. Trazabilidad:} el sistema debe exponer las plantillas
    generadas antes de ejecutar el despliegue, de forma que el usuario pueda
    revisarlas y auditarlas.

    \item \textbf{RNF-05. Mantenibilidad:} la arquitectura interna del sistema
    debe facilitar la incorporación de nuevos tipos de recursos y la corrección
    de equivalencias entre proveedores, sin requerir cambios transversales en el
    código. Cada proveedor y cada tipo de recurso deben poder modificarse de forma
    independiente.
\end{itemize}

\section{Líneas de evolución futura}

Los siguientes aspectos quedan explícitamente fuera del alcance de este trabajo
y se identifican como líneas de evolución futura:

\begin{itemize}
    \item Soporte a proveedores adicionales: \textbf{Google Cloud Platform},
    Oracle Cloud, Alibaba Cloud y otros.

    \item Generación de Manifest desde \texttt{docker-compose.yml} u otras
    fuentes de definición de entornos locales.

    \item Interfaz web con gestión visual de infraestructura, diagramas
    de arquitectura interactivos y paneles de monitorización.

    \item Componente agéntico: asistente en lenguaje natural, integración
    con agentes externos mediante MCP y sugerencias de optimización
    basadas en el contexto del despliegue.

    \item Monitorización del gasto cloud y detección de cloud sprawl en
    infraestructura ya desplegada.
    
    \item Gestión de estado propia de CloudPorter: seguimiento de manifests
    desplegados, proveedor y versión, con almacén embebido o backend remoto
    para entornos multi-usuario.

    \item Cobertura de lock-in de código: integración con herramientas de
    abstracción a nivel de SDK para cubrir servicios con dependencia a nivel
    de código.

\end{itemize}
